% Guia de construção — reconstruir o sistema de raiz, por fases, com o código real.
%
% ⚠️ POR QUE É QUE ESTE DOCUMENTO EXISTE, e em que difere de `tese/guia/`.
% Aquele ensina a DEFENDER: os números de cor, as perguntas do júri, o que não dizer. Este ensina
% a CONSTRUIR. São complementares, e a razão de existirem os dois é que as perguntas que um
% arguente faz sobre COMO uma coisa foi feita não se respondem com um número decorado.
%
% ⚠️ TODO O CÓDIGO AQUI É VERBATIM. Foi copiado dos ficheiros e conferido por máquina, duas vezes:
% uma a exigir que cada linha exista no ficheiro, outra a exigir que os blocos sejam CONTÍGUOS e
% pela ordem certa. A segunda passagem apanhou um excerto que colava duas funções e saltava uma
% terceira sem marcar o corte — e a primeira tinha-o dado por bom. Onde há um corte, ele está
% marcado com `# ...` e nunca escondido.
%
% Regras do projecto que valem aqui: PT-PT, decimais com PONTO, zero travessões na prosa.

\documentclass[aspectratio=169]{beamer}
\usetheme{Madrid}
\usecolortheme{seahorse}
\usepackage[T1]{fontenc}
\usepackage[portuguese,shorthands=off]{babel}
\usepackage{booktabs}
\usepackage{listings}
\usepackage{xcolor}
\usepackage{tikz}
\graphicspath{{../figures/}{../slides/}}

\definecolor{marca}{HTML}{0A8F52}
\newcommand{\gator}{Investi\textcolor{marca}{G}ator}
\setbeamertemplate{navigation symbols}{}
% Um título de frame a `\large` custa cerca de meia dúzia de pontos de altura em todos os slides,
% e estes são densos por natureza: cada um tem um excerto de código inteiro. Num guia de leitura
% o título não precisa de ser grande, e a folga que isto liberta é o que faz o "Verifica-se com"
% caber acima do rodapé em vez de lhe assentar em cima.
\setbeamerfont{frametitle}{size=\small}

% O mesmo estilo de código da dissertação, para o aluno reconhecer os excertos quando os
% reencontrar lá. `literate` é obrigatório: sem ele os acentos dos comentários reais saem partidos.
\lstset{
  basicstyle=\ttfamily\tiny,
  keywordstyle=\bfseries,
  commentstyle=\itshape\color{black!55},
  stringstyle=\color{black!70},
  frame=single, rulecolor=\color{black!25}, framesep=3pt,
  backgroundcolor=\color{black!3},
  % Os excertos passaram a ser cortados do ficheiro em vez de escritos a mão, e ficaram mais
  % altos: cinco frames transbordavam, o pior em 22pt, o que empurra o "Verifica-se com" para
  % cima do rodapé. O espaço à volta do bloco é onde há folga sem mexer no código.
  aboveskip=1mm, belowskip=1mm,
  showstringspaces=false, breaklines=true, columns=fullflexible,
  language=Python, literate=%
    {á}{{\'a}}1 {ã}{{\~a}}1 {ç}{{\c c}}1 {é}{{\'e}}1 {í}{{\'i}}1
    {ó}{{\'o}}1 {õ}{{\~o}}1 {ú}{{\'u}}1 {â}{{\^a}}1 {ê}{{\^e}}1
    {à}{{\`a}}1 {ô}{{\^o}}1 {É}{{\'E}}1 {Ó}{{\'O}}1
    % ⚠️ Estes dois estão em código REAL e por isso ficam. O ponto médio é do docstring do
    % calibrador (`a·score + b`) e o travessão é do docstring do `assign_splits`. A regra do
    % projecto de zero travessões aplica-se à PROSA, não a texto citado: reescrever o comentário
    % de um ficheiro para o pôr num guia seria deixar de ser verbatim, que é a única coisa que
    % este documento promete.
    % Idem: a seta esta numa mensagem de registo real do prices.py, e o "menor ou igual" num
    % docstring do dataset.py. Substitui-los por "->" e "<=" seria reescrever o ficheiro.
    {·}{{\textperiodcentered}}1 {—}{{---}}1
    {→}{{$\rightarrow$}}1 {≤}{{$\leq$}}1
}

% ⚠️ `\fonte` não é decoração. Faz duas coisas: diz ao leitor que ficheiro abrir para ver o
% excerto no seu contexto, e é o que permite ao `scripts/check_guia_codigo.py` conferir, a cada
% compilação, que o bloco seguinte continua a existir VERBATIM nesse ficheiro. Sem a marca, a
% promessa deste guia (código real, não reescrito) seria só uma afirmação minha.
% O `\detokenize` é obrigatório e não é preciosismo: os caminhos têm `_` (market_data,
% news_fetcher, onnx_embedder) e um `_` fora de modo matemático é um erro de LaTeX. Escapá-los à
% mão em cada chamada funcionaria, mas partiria o `check_guia_codigo.py`, que lê estes caminhos
% para ir ao ficheiro: passaria a procurar `market\_data`, que não existe.
\newcommand{\fonte}[1]{\vspace{0.5mm}{\tiny\ttfamily\color{black!55}\detokenize{#1}}\vspace{-1mm}}
\newcommand{\linhaq}[1]{\vspace{1mm}{\scriptsize\textbf{A linha que interessa.} #1}}
% O `\par` põe isto em linha própria sem gastar espaço vertical: sem ele o comando ficava colado
% ao fim da frase anterior e lia-se como parte dela. Um `\vspace` resolveria a leitura e custaria
% meio milímetro em cada frame, que é exactamente o que separa dois deles de transbordarem.
\newcommand{\verifica}[1]{{\scriptsize\par\textbf{Verifica-se com.} \texttt{#1}}}

\title[Guia de construção]{Reconstruir o \gator{} de raiz}
\subtitle{\normalsize Dez fases, com o código que está mesmo lá}
\author[H.\ Santos]{Henrique José da Silva Santos}
\institute[ISEP]{MEIA, Instituto Superior de Engenharia do Porto}
\date{\footnotesize Guia de estudo}

\begin{document}

\begin{frame}\titlepage\end{frame}

% ══════════════════════════════════════════════════════════════════════════════
\begin{frame}{Para que serve este guia, e para que não serve}
\small
O outro guia ensina-te a \textbf{defender}: os números de cor, as perguntas prováveis, o que não
dizer. Este ensina-te a \textbf{construir}.
\vspace{2mm}

\begin{center}
\textbf{A diferença importa numa prova.}\\
Uma pergunta sobre \emph{quanto deu} responde-se com um número.\\
Uma pergunta sobre \emph{como fizeste} não.
\end{center}
\vspace{3mm}

São dez fases, pela ordem em que o sistema foi realmente construído. Cada uma traz:
\vspace{1mm}
\begin{itemize}\small
  \item o \textbf{problema} que ela resolve, numa frase;
  \item o \textbf{código real}, copiado dos ficheiros e não reescrito;
  \item \textbf{a linha que interessa}, ou seja onde é que a garantia é mesmo feita;
  \item \textbf{como se verifica}, com um comando que podes correr.
\end{itemize}
\vspace{2mm}

{\footnotesize Os erros que cada fase custou estão reunidos no penúltimo \emph{slide}, de
propósito: lidos de seguida, são a parte deste trabalho que mais se aprende.}
\end{frame}

% ══════════════════════════════════════════════════════════════════════════════
\begin{frame}[fragile]{Fase 1: de preços para retornos}
{\footnotesize \$500 é muito numa empresa e pouco noutra: o sistema trabalha em
\textbf{retornos}, e tenta cinco fontes gratuitas por ordem fixa.}

\fonte{investigator/market_data/prices.py}
\begin{lstlisting}
def get_price_history(ticker: str, period: str = "6mo", interval: str = "1d") -> pd.DataFrame:
# ...
    try:
        df = _yf_history(ticker, period=period, interval=interval)
        _LAST_SOURCE[ticker.upper()] = "yfinance"
        return df
    except Exception as exc:  # noqa: BLE001
        if interval != "1d":
            raise RuntimeError(f"Sem dados de preços para o ticker '{ticker}'.") from exc
# ...
        df, fonte = fallback_daily(ticker, start, end)
        _LAST_SOURCE[ticker.upper()] = fonte
        print(f"[precos {ticker}] yfinance indisponível → servido por {fonte} ({len(df)} dias)")
        return df
# ...
def log_returns(close: pd.Series) -> pd.Series:
    """Retornos logarítmicos a partir da série de preços de fecho."""
    return np.log(close / close.shift(1)).dropna()
\end{lstlisting}

\linhaq{O \texttt{shift(1)}. Cada retorno divide o fecho de hoje pelo do dia \textbf{anterior} e
nunca por um posterior: nenhum valor futuro entra na norma que a fase seguinte constrói.}
\verifica{pytest tests/test\_prices\_fallback.py}
\end{frame}

% ══════════════════════════════════════════════════════════════════════════════
\begin{frame}[fragile]{Fase 2: o que conta como movimento invulgar}
{\footnotesize $2\%$ é banal na Tesla e enorme na Johnson \& Johnson, portanto um limiar fixo em
percentagem mediria \emph{volatilidade} e não \emph{raridade}. Mede-se em desvios-padrão face aos
últimos vinte dias da própria empresa.}
\vspace{1mm}

\fonte{investigator/anomaly_detector/detector.py}
\begin{lstlisting}
def _score(last: float, mu: float, sigma: float, threshold: float) -> tuple[float, bool, bool]:
# ...
    if sigma > 0.0:
        z = (last - mu) / sigma
        return z, abs(z) > threshold, False
    moved = not math.isclose(last, mu, rel_tol=0.0, abs_tol=1e-15)
    return 0.0, moved, True
# ...
    recent = r.iloc[-window - 1 : -1]  # janela ANTES do último (sem lookahead)
    last = float(r.iloc[-1])
    mu = float(recent.mean())
    sigma = float(recent.std(ddof=1))
    z, is_anomaly, zero_variance = _score(last, mu, sigma, threshold)
# ...
        zero_variance=zero_variance,
\end{lstlisting}

\linhaq{O \texttt{-1} do fim da fatia \texttt{r.iloc[-window - 1 : -1]}, que \textbf{exclui o
próprio dia julgado}. Com \texttt{r.iloc[-window:]} o salto entrava no cálculo da sua própria
régua, inflacionava o \texttt{sigma} e encolhia o seu próprio $z$. Se a norma for constante,
repetir o valor não dispara; afastar-se dele dispara, mas o produto diz que o $z$ é indefinido.}
\verifica{pytest tests/test\_anomaly\_detector.py}
\end{frame}

% ══════════════════════════════════════════════════════════════════════════════
\begin{frame}[fragile]{Fase 3: deitar fora o que não é da empresa}
{\footnotesize A fonte etiqueta mal. Chegavam notícias de escritórios de advogados etiquetadas
como \emph{AMD}, e listas de \emph{top movers} repetidas para vários tickers. Sem este filtro,
$67\%$ do que entra é ruído, e um precedente construído sobre ruído é pior do que nenhum.}
\vspace{1mm}

\fonte{investigator/news_fetcher/relevance.py}
\begin{lstlisting}
def is_relevant(headline: str, ticker: str) -> bool:
    """True se a manchete é plausivelmente SOBRE esta empresa (e não boilerplate de mercado)."""
    if not headline or not headline.strip():
        return False
    low = headline.lower()
    if _BOILERPLATE_RE.search(low):
        return False
    terms = [ticker.lower()] + COMPANY_NAMES.get(ticker.upper(), [])
    return any(_mentions(low, t) for t in terms)
\end{lstlisting}

\linhaq{A ordem. O \texttt{\_BOILERPLATE\_RE} corre \textbf{antes} da procura de menções, e tem
de correr: uma lista de \emph{top movers} menciona a empresa, portanto passaria a última linha
sem dificuldade. Trocar as duas deixaria entrar exactamente o ruído que este filtro existe para
apanhar, e nada avisaria. Repare-se ainda no \texttt{COMPANY\_NAMES}: procurar só o \emph{ticker}
perderia toda a notícia que escreve \emph{Nvidia} e nunca \emph{NVDA}.}
\verifica{pytest tests/test\_relevance.py}
\end{frame}

% ══════════════════════════════════════════════════════════════════════════════
\begin{frame}[fragile]{Fase 4: uma frase como um ponto no espaço}
{\footnotesize Duas notícias sobre o mesmo assunto raramente partilham palavras. Um codificador de
frases já treinado transforma cada título em $384$ números, de forma a que proximidade entre pontos
seja proximidade de \emph{significado}.}
\vspace{1mm}

\fonte{investigator/historical_kb/onnx_embedder.py}
\begin{lstlisting}
def masked_mean_pool(hidden: np.ndarray, attention_mask: np.ndarray) -> np.ndarray:
    """Mean pooling com máscara + normalização L2 — igual ao sentence-transformers.

    hidden: (n, seq, dim); attention_mask: (n, seq) com 1 nos tokens reais. Puro numpy,
    testável offline (tests/test_onnx_embedder.py cobre a matemática sem o modelo).
    """
    mask = attention_mask[..., None].astype("float64")
    summed = (hidden.astype("float64") * mask).sum(axis=1)
    counts = np.clip(mask.sum(axis=1), 1e-9, None)
    pooled = summed / counts
    norms = np.linalg.norm(pooled, axis=1, keepdims=True)
    return pooled / np.clip(norms, 1e-12, None)
\end{lstlisting}

\linhaq{A última linha, a normalização. Depois dela todos os vectores têm comprimento $1$, e por
isso o cosseno entre dois passa a ser um simples produto interno. Não é uma optimização: é o que
faz com que \emph{semelhança} deixe de depender do comprimento do título.}
\verifica{pytest tests/test\_onnx\_embedder.py}
\end{frame}

% ══════════════════════════════════════════════════════════════════════════════
\begin{frame}[fragile]{Fase 5: encontrar o precedente}
{\footnotesize Com cada título já como um ponto, \emph{parecido} passa a ser uma conta: o cosseno
do ângulo entre dois vectores. Um resultado perto de $1$ quer dizer que apontam para o mesmo sítio.}
\vspace{1mm}

\fonte{investigator/correlation_engine/similarity.py}
\begin{lstlisting}
def cosine_similarity(a, b) -> float:
    """Similaridade do cosseno entre dois vetores 1D. Vetor nulo → 0.0 (sem direção)."""
    a = np.asarray(a, dtype="float64")
    b = np.asarray(b, dtype="float64")
    na = float(np.linalg.norm(a))
    nb = float(np.linalg.norm(b))
    if na == 0.0 or nb == 0.0:
        return 0.0
    return float(np.dot(a, b) / (na * nb))
\end{lstlisting}

\linhaq{A guarda \texttt{if na == 0{,}0 or nb == 0{,}0}. Um vector nulo não tem direcção, logo o
ângulo não existe e a divisão rebentaria. Devolver $0{,}0$ diz \emph{não se parece com nada}, que é
a leitura correcta, e mantém a função total: nenhuma consulta pode derrubar o sistema por causa de
um registo degenerado.}
\verifica{pytest tests/test\_similarity.py}
\end{frame}

% ══════════════════════════════════════════════════════════════════════════════
\begin{frame}[fragile]{Fase 6: medir o que aconteceu depois}
{\footnotesize Esta é a fase que separa este trabalho de uma previsão. Mede-se o que o preço
\textbf{fez} nos $1$, $3$ e $5$ dias a seguir a uma notícia \emph{passada}. É observação
verificável, não projecção.}
\vspace{1mm}

\fonte{investigator/correlation_engine/event_study.py}
\begin{lstlisting}
def post_event_returns(close: pd.Series, event_idx: int,
                       horizons: tuple[int, ...] = (1, 3, 5)) -> dict[int, float]:
    # ...
    p0 = float(close.iloc[event_idx])
    out: dict[int, float] = {}
    for h in horizons:
        j = event_idx + h
        out[h] = float(close.iloc[j] / p0 - 1.0) if 0 <= j < len(close) else float("nan")
    return out
# ...
def abnormal_returns(close: pd.Series, market_close: pd.Series, event_idx: int,
                     horizons: tuple[int, ...] = (1, 3, 5)) -> dict[int, float]:
    # ...
    own = post_event_returns(close, event_idx, horizons)
    mkt = post_event_returns(market_close, event_idx, horizons)
    return {h: own[h] - mkt[h] for h in horizons}
\end{lstlisting}

\linhaq{O \texttt{float("nan")} quando \texttt{j} sai da série. Uma notícia dos últimos cinco dias
\textbf{ainda não tem} desfecho a $+5$; devolver zero, ou o último preço disponível, inventaria um
resultado. O \texttt{NaN} propaga-se e obriga quem consome a decidir o que fazer com a ausência.}
\verifica{pytest tests/test\_event\_study.py}
\end{frame}

% ══════════════════════════════════════════════════════════════════════════════
\begin{frame}[fragile]{Fase 7: foi a empresa, ou foi o mercado?}
{\footnotesize Um movimento reparte-se em três: o que veio do mercado inteiro, o que veio do setor,
e o que sobra e é da própria empresa. Os betas são estimados só com dias \textbf{anteriores}.}
\vspace{1mm}

\fonte{investigator/correlation_engine/decomposition.py}
\begin{lstlisting}
def _shrink(beta_raw: float, std_err: float, prior: float,
            prior_sd: float = PRIOR_BETA_SD) -> float:
    """Encolhe `beta_raw` na direção de `prior`, pesando pela precisão da estimativa."""
    if not np.isfinite(std_err) or std_err <= 0:
        return beta_raw  # ajuste exato: nada a encolher
    w = prior_sd**2 / (prior_sd**2 + std_err**2)
    return w * beta_raw + (1.0 - w) * prior
# ...
    total = float(t[-1])
    r_m_today = float(m[-1])
    r_s_today = float(s[-1]) if s is not None else 0.0

    # Janela de estimação: estritamente ANTES do dia explicado (anti-lookahead).
    hist_t, hist_m = t[:-1][-window:], m[:-1][-window:]
    hist_s = s[:-1][-window:] if s is not None else None
\end{lstlisting}

\linhaq{O peso \texttt{w} do \texttt{\_shrink}. Um beta estimado em vinte dias é ruidoso, e numa
janela agitada a AMD deu $\beta = 4{,}43$. Cortar em $\pm 4$ seria outra constante por justificar;
aqui o encolhimento é \textbf{proporcional à precisão da estimativa}: com erro-padrão pequeno o
beta fica intacto, com erro grande aproxima-se da referência.}
\verifica{pytest tests/test\_decomposition.py}
\end{frame}

% ══════════════════════════════════════════════════════════════════════════════
\begin{frame}[fragile]{Fase 8: um conjunto de treino sem o futuro lá dentro}
{\footnotesize É aqui que os trabalhos com dados temporais se enganam, e o engano não dá erro:
dá um resultado bom de mais.}

\fonte{investigator/triage/dataset.py}
\begin{lstlisting}
def assign_splits(dates: pd.Series, train_frac: float = 0.70, val_frac: float = 0.15,
                  embargo_days: int = 5) -> pd.Series:
    """Divisão TEMPORAL por dias únicos (70/15/15) com embargo entre blocos.

    Todas as linhas do mesmo dia ficam no mesmo bloco (evita fuga por notícias do mesmo dia).
    O embargo remove `embargo_days` dias únicos APÓS cada fronteira (rotulados "embargo",
    excluídos do treino e da avaliação) — protege contra sobreposição das janelas de rótulo
    (h≤5) entre blocos.
    """
    d = pd.to_datetime(pd.Series(dates).reset_index(drop=True))
    unique_days = np.array(sorted(d.unique()))
    n = len(unique_days)
    i1 = int(n * train_frac)
    i2 = int(n * (train_frac + val_frac))
    train_days = set(unique_days[:i1])
    emb1 = set(unique_days[i1 : i1 + embargo_days])
    val_days = set(unique_days[i1 + embargo_days : i2])
    emb2 = set(unique_days[i2 : i2 + embargo_days])
\end{lstlisting}

\linhaq{\texttt{unique\_days}, e não as linhas. Dividir por linhas poria duas notícias
\textbf{do mesmo dia}, com o mesmo rótulo, uma no treino e outra no teste: o modelo veria a
resposta. O embargo trata do segundo problema, o rótulo olhar cinco dias em frente.}
\verifica{pytest tests/test\_triage\_dataset.py}
\end{frame}

% ══════════════════════════════════════════════════════════════════════════════
\begin{frame}[fragile]{Fase 9: treinar, e depois calibrar}
{\footnotesize Um modelo dá uma pontuação, e uma pontuação não é uma probabilidade. Para o alerta
poder dizer \emph{$39\%$} é preciso um segundo ajuste, feito num bloco reservado que o modelo
nunca viu.}
\vspace{1mm}

\fonte{investigator/triage/model.py}
\begin{lstlisting}
class PlattCalibrator:
    """p_calibrada = sigmóide(a·score + b), com (a,b) ajustados na validação."""

    a: float
    b: float

    def __call__(self, scores: np.ndarray) -> np.ndarray:
        z = self.a * np.asarray(scores, dtype="float64") + self.b
        return 1.0 / (1.0 + np.exp(-z))


def fit_platt(scores_val: np.ndarray, y_val: np.ndarray, seed: int = 42) -> PlattCalibrator:
    lr = LogisticRegression(max_iter=1000, random_state=seed)
    lr.fit(np.asarray(scores_val, dtype="float64").reshape(-1, 1), y_val)
    return PlattCalibrator(a=float(lr.coef_[0][0]), b=float(lr.intercept_[0]))
\end{lstlisting}

\linhaq{O \texttt{scores\_val}, no nome do argumento. A calibração ajusta-se na
\textbf{validação} e nunca no treino: ajustada no treino, corrigiria o sobreajustamento em vez da
escala, e as probabilidades sairiam confiantes de mais onde o modelo está errado. Aqui
$a = 3{,}700$ e $b = -2{,}313$.}
\verifica{pytest tests/test\_triage\_model.py}
\end{frame}

% ══════════════════════════════════════════════════════════════════════════════
\begin{frame}[fragile]{Fase 10: observar o que corre, e deixar-se contrariar}
{\footnotesize Sem esta fase o trabalho acabava numa métrica de teste. Cada decisão real é
registada, e dias depois confrontada com o que o preço \emph{fez}. Foi assim que se soube que a
porta \textbf{não estava a ajudar}.}
\vspace{1mm}

\fonte{investigator/triage/postval.py}
\begin{lstlisting}
    rep: dict = {
        "n": n,
        "n_mantidas": len(kept),
        "precisao_mantidas": (float(np.mean([d["label"] for d in kept])) if kept else float("nan")),
        "n_suprimidas": len(dropped),
        "precisao_suprimidas": (float(np.mean([d["label"] for d in dropped]))
                                if dropped else float("nan")),
        "base_rate": (float(np.mean([d["label"] for d in labeled])) if labeled else float("nan")),
        "brier": (float(np.mean([(d["prob"] - d["label"]) ** 2 for d in with_prob]))
                  if with_prob else float("nan")),
        "calibracao": [],
    }
\end{lstlisting}

\linhaq{A \texttt{precisao\_suprimidas}. Comparar as mantidas com a taxa-base responde a
\emph{a porta acerta?}, que é a pergunta fácil. A pergunta que interessa é \emph{a porta escolhe
melhor do que aquilo que deitou fora?}, e essa exige medir também as suprimidas. Medido: mantidas
$0{,}589$ contra suprimidas $\mathbf{0{,}617}$, ou seja o que a porta rejeitou era \emph{mais}
material.}
\verifica{python scripts/post\_validate.py}
\end{frame}

% ══════════════════════════════════════════════════════════════════════════════
% ⚠️ ESTA TABELA ESTAVA NUM SÓ FRAME E CORTAVA AS DUAS ÚLTIMAS LINHAS, incluindo a dos dezanove
% dias, que é a mais forte de todas. O aviso era `Overfull \vbox 70pt`, ou seja não é um erro e
% não para a compilação: só se vê a renderizar a página. Dividida em duas, quatro linhas cada.
\begin{frame}{Os erros que cada fase custou (1/2)}
\small
Lidos de seguida, são a parte deste trabalho que mais se aprende, e são a melhor resposta a
\emph{``o que é que correu mal?''}. Nenhum deles deu erro nenhum.
\vspace{3mm}

{\scriptsize
\begin{tabular}{@{}p{0.15\textwidth}p{0.80\textwidth}@{}}
\toprule
\textbf{fase} & \textbf{o erro, e porque é que passou despercebido} \\
\midrule
1. preços & A última fonte da cadeia devolvia um resultado \emph{vazio} como sucesso. Preços
vazios propagavam-se em silêncio, sem uma única excepção. Numa cadeia de recurso, a resposta
vazia tem de ser tratada como falha. \\
\addlinespace
2. detetor & A Microsoft a $+4{,}82\%$ com o cartão a dizer \emph{``an ordinary day''}. O $z$ era
$+1{,}11$, abaixo do limiar, mas só $5$ dos últimos $249$ dias se tinham movido tanto. Duas réguas
a discordar, e a interface escolhia em silêncio a mais tranquilizadora. \\
\addlinespace
3. relevância & Sem o filtro, a mesma história contava duas vezes e um escritório de advogados
era notícia da AMD. \\
\addlinespace
5. precedentes & A deduplicação era de \emph{texto exacto}: a mesma história escrita por dois
meios continuava a contar como duas observações, e o alerta afirmava mais evidência do que
tinha. \\
\bottomrule
\end{tabular}}
\end{frame}

\begin{frame}{Os erros que cada fase custou (2/2)}
\vspace{2mm}
{\scriptsize
\begin{tabular}{@{}p{0.15\textwidth}p{0.80\textwidth}@{}}
\toprule
\textbf{fase} & \textbf{o erro, e porque é que passou despercebido} \\
\midrule
6. impacto & O impacto é medido por (empresa, dia), portanto três títulos do mesmo dia partilham
o mesmo valor \emph{por construção}. Em $36{,}8\%$ dos alertas os casos assentavam em menos dias
distintos do que casos mostrados: o alerta afirmava mais evidência independente do que tinha. \\
\addlinespace
7. decomposição & O motor era escolhido pela maior componente \textbf{em módulo}, e uma empresa
a subir com o setor a puxar ao contrário saía descrita como \emph{``foi o setor''}. O sinal
importa, e o módulo deita-o fora. \\
\addlinespace
9. calibração & O ajuste é feito numa validação a $47{,}0\%$ de prevalência e aplicado a um teste a
$37{,}8\%$. A ordenação não muda, mas os limiares implantados assentam nessa escala, e por isso a
probabilidade média prevista sai $0{,}050$ acima da observada. \\
\addlinespace
10. observação & Os dois registos que alimentam esta fase eram escritos em disco local, num
contentor cujo disco é efémero e pertence a outro processo. Estiveram parados
\textbf{dezanove dias sem um único erro registado}, e foi encontrado a olhar, não por um alarme. \\
\bottomrule
\end{tabular}}
\vspace{3mm}

\begin{center}\footnotesize
O padrão é sempre o mesmo: cada peça cumpria o seu contrato, e o que estava errado era uma
suposição sobre a ligação entre elas que ninguém tinha escrito em lado nenhum.
\end{center}
\end{frame}

% ══════════════════════════════════════════════════════════════════════════════
\begin{frame}{Se abrirem o repositório: o que lá está e a tese não reivindica}
\small
É uma pergunta possível, e a resposta honesta é melhor do que a evasiva.
\vspace{3mm}

Existe em \texttt{investigator/intelligence/} e \texttt{investigator/narrator/} uma camada que usa
um modelo de linguagem para escrever prosa, com uma verificação que rejeita qualquer frase cujos
números não estejam ancorados num facto citado. Foi construída, foi atacada por um exercício
adversário, e \textbf{$23$ de $23$ ataques foram bloqueados}.
\vspace{3mm}

\begin{center}
\textbf{E está desligada em produção, de propósito.}
\end{center}
\vspace{3mm}

Duas razões, e as duas estão no documento. A primeira é de posição: o trabalho defende que uma
frase gerada é uma \emph{afirmação}, e este sistema entrega afirmações com a evidência ao lado.
A segunda é de honestidade: a dissertação \textbf{não reivindica} essa camada em lado nenhum, e
manter no ar aquilo que o documento não descreve seria dívida.
\vspace{3mm}

{\footnotesize A frase a dizer: \emph{``está lá, funciona, e decidi não a expor porque contradiz a
posição que o trabalho assume. As rotas foram retiradas e a configuração está a falso.''}}
\end{frame}

% ══════════════════════════════════════════════════════════════════════════════
\begin{frame}{O fio que liga as dez fases}
\small
Se te esquecerem de tudo o resto, é isto que a construção ensina.
\vspace{3mm}

\begin{enumerate}
  \item Cada fase tem \textbf{uma linha} onde a garantia é feita, e quase sempre é uma fatia, uma
        guarda ou a escolha de que dados entram. Não é o algoritmo.
  \item Os erros que custaram mais tempo \textbf{não davam erro nenhum}. Davam um resultado
        plausível.
  \item A parte difícil não foi escolher a técnica. Foi montar a avaliação que
        \textbf{consegue dizer que não}.
\end{enumerate}
\vspace{4mm}

\begin{center}
\footnotesize Guia irmão: \texttt{tese/guia/} ensina a defender.\\
Este ensina a construir. Levas os dois.
\end{center}
\end{frame}

\end{document}
